\section{Two-class and prescribed-cross consequences}\label{sec:results}
Fix integers $a\ge0$, $b\ge1$, and $D=a+b\ge3$, and put
$\delta=a-b$, $n=2N$. There are two planted classes of size $N$;
their independent internal factors are uniform simple $a$-regular,
and the independent cross factor is uniform simple $b$-regular
bipartite. Degree-zero factors are absent. All limits are through
admissible sizes, with $aN$ even and the simple factors existing.
Before hiding the plant write $\sigma=(\one_N,-\one_N)$.
Deterministically,
\begin{equation}\label{eq:planted}
 A\one_n=D\one_n,\qquad A\sigma=\delta\sigma.
\end{equation}
A sign vector $\tau$ satisfies $A\tau=\delta\tau$ exactly when
each vertex has $a$ neighbors of its own sign and $b$ of the other.
It is automatically balanced, since $(D-\delta)\one^T\tau=2b\one^T\tau=0$.

\begin{theorem}[All two-class parameters]\label{thm:main-recovery}
For every fixed $a\ge0,b\ge1,D\ge3$, all valid unordered $(a,b)$
bisections can be listed by an always-correct Las Vegas algorithm in
unconditional expected $\widetilde O_{a,b}(n^2)$ bit work and
expected $\widetilde O_{a,b}(n)$ workspace. Returning any member
recovers the plant with probability $1-O_{a,b}(N^{-\nu})$, where
\[
 \nu=(D-2)m,\qquad
 m=\begin{cases}a-b+1,&a\ge b,\\b-a,&a<b.\end{cases}
\]
Under a uniformly hidden balanced plant this selection attains the
optimal Bayes risk exactly. A deterministic polynomial estimator with
bounded $O_{a,b}(n^3+n^{\nu+2})$ elementary work and $O(n^2)$ storage
attains its leading constant. All admissible sizes are included.
No assortativity, real-corank, or adjacency outlier assumption is needed.
The unknown pair $a,b$ can also be recovered with high probability
in expected nearly quadratic work.
\end{theorem}
\begin{proof}[Algorithmic implications]
Theorem~\ref{thm:frame-bounded-moment} with scalar target $\delta$
bounds the actual-law finite-field kernel moment. Apply
Corollary~\ref{fk:cor:moment}, check the original integer sign
equation, and identify opposite signs. The sharp ambiguity theorem
below gives the recovery rate, uniform posterior and Bayes assertion;
Section~\ref{sec:enumeration} supplies the deterministic cap.
For unknown parameters, read $D$ from the graph and enumerate the
finite family $a\ge0,b\ge1,a+b=D$. Alternatively apply the unknown-frame
algorithm with class bound two. On its rigidity event, every nonempty
two-class candidate is the plant and its integer neighbor counts fix
$a,b$. The finite-family moment theorem preserves the expected work.
\end{proof}

\begin{theorem}[Bounded scalar moment]\label{thm:general-moment}
For every fixed $a\ge0,b\ge1,D\ge3$ and prime $p$, in both the
configuration and independent uniform-simple laws,
\begin{equation}\label{eq:poly-moment}
 \sup_N\E|\ker_{\Fp}(A-\delta I)|\le C_{a,b,p}<\infty.
\end{equation}
The simple-law supremum ranges over sizes where its factors exist.
\end{theorem}
This is the two-equal-class specialization of
Theorem~\ref{thm:frame-bounded-moment}. The Boolean rigidity theorem
below proves failure $O(N^{-\nu})$ and its sharp form computes
the exact leading constant. Neither statement asserts real
eigenspace simplicity.

For the separate prescribed-cross spectral theorem, fix a connected
symmetric $Q=(d_{ij})\in\Z_{\ge0}^{k\times k}$ with $d_{ii}\ge3$.
Its off-diagonal support is connected, with the one-type convention.
There are $k$ equal classes of size $N$, with $d_{ii}N$ even.
Internal blocks are independent uniform simple $d_{ii}$-regular.
Every cross block is an arbitrary prescribed simple $d_{ij}$-regular
bipartite graph independent of all internal randomness. Set
\begin{equation}\label{eq:J}
 Jc=(c_1\one_N,\ldots,c_k\one_N),\qquad W=\im J.
\end{equation}
Then $AJ=JQ$; the total graph is regular precisely when $Q$ has
constant row sum.

\begin{theorem}[Uniform fixed-target exclusion]\label{thm:frame}
For every fixed real $\lambda$, with
$\ell=\dim_\R\ker(Q-\lambda I)$,
\begin{equation}\label{eq:frame-main}
 \sup_{(B_{ij})}\Prob\{
 \ker_\R(A-\lambda I)\ne J\ker_\R(Q-\lambda I)\}\longrightarrow0
\end{equation}
along every admissible sequence. All quotient eigenspaces thus have
exactly their old multiplicities simultaneously with high probability.
Cross blocks may instead have any joint law independent of the
internal blocks.
\end{theorem}

\begin{corollary}[Known-frame rational recovery]\label{cor:frame-recovery}
Given $A,Q$ in Theorem~\ref{thm:frame}, exact rational computation
of $\ker_\Q m_Q(A)$ and grouping its equal rows recovers the planted
classes with high probability, uniformly over the cross blocks.
Here $m_Q$ is the squarefree minimal polynomial of $Q$.
Repeated and irrational quotient eigenvalues are allowed.
\end{corollary}

\begin{corollary}[Two-class prescribed cross graph]\label{cor:two-frame}
For fixed $a\ge3,b\ge1$, and arbitrary prescribed simple $b$-regular
cross block independent of its internal graphs,
\[
 \sup_B\Prob\{\ker_\R(A-(a-b)I)\ne\R\sigma\}\longrightarrow0.
\]
This also holds for $(a,b)=(2,1)$, hence for every $a>b>0$.
Every admissible size is covered. Rational nullspace computation
recovers the plant with high probability. The $a\ge3$ statement
itself requires no assortativity.
\end{corollary}

The prescribed-cross proof uses direct all-target internal configuration
mixing and positive-probability simplicity conditioning, including odd
admissible $N$. The cubic case uses the sharp regular-frame scalar
moment and the cross-matching gauge. This robust quantifier is distinct
from the all-parameter random-cross theorem; it is not asserted for
every prescribed cross graph when $a=0$ or $a=1$.

